\documentclass[ddcfooter,nototalpage]{tudbeamer}
\usepackage{german}
\usepackage[utf8]{inputenc}
\usepackage[T1]{fontenc}
\usepackage{tikz}

\begin{document}
\einrichtung{Fakultät Mathematik und Naturwissenschaften}
\title{SCRUM - Tag 6}
\author{Team 29: Cherno Alpha \\
Friedrich Zahn \\
David Tucholski \\
Thomas Adlmaier \\
Sebastian Schmidt
}
\date{23.3.2015}
\maketitle

\section{Aktueller Stand}
\begin{frame}
\frametitle{Aktueller Stand}

\textbf{Roboterteam:}
\begin{itemize}
\item Sämtliche Bewegungsfunktionen fertiggestellt
\item Kalibrierungsfunktionen für S/W-Erkennung und Umdrehungszahlen
\item Roboter kann jetzt alle an ihn gestellten Aufgaben erfüllen
\item Code aufgeräumt
\end{itemize}\ \linebreak

\textbf{Algorithmusteam:}
\begin{itemize}
\item Algorithmus komplett
\item löst Labyrinth eindeutig und erfüllt die Aufgaben 3 und 5
\item Code aufgeräumt
\end{itemize}
\end{frame}

\begin{frame}
\normalsize
\textbf{Code-Heirat}
\begin{itemize}
\item steht noch am Anfang
\item Roboter kommt bisher nur zum ersten Knoten
\end{itemize}

\vspace{.5cm}
\begin{tikzpicture}[scale=.8]
\draw [yellow, line width = .2cm] (0,0) circle[radius=.5cm];
\draw [yellow, line width = .2cm] (.5,0) circle[radius=.5cm];
\end{tikzpicture}

\begin{itemize}
\item Abschlussbericht angefangen
\end{itemize}
\end{frame}

\section{Tagesziele}
\begin{frame}
\frametitle{Tagesziele}
\normalsize

\begin{itemize}
\item Weiterhin versuchen den Algorithmus auf den Roboter zu laden, Fehlersuche
\item Code weiterhin vereinfachen beim Zusammenführen
\item Abschlussbericht weiterführen
\end{itemize}
\end{frame}

\section{Probleme}
\begin{frame}
\frametitle{Probleme}
\normalsize

\textbf{Roboterteam}
\begin{itemize}
\item unbekannter Fehler wenn beide Codes vereint werden
\end{itemize}
\end{frame}
\end{document}